\begin{lemma}[Reduction to the top homogeneous component]
Let $n\in\mathbb{N}$, let $p\in\mathbb{R}$ and let
$P\in\mathbb{R}[X_1,\dots,X_n]$ with $\deg P\ge2$, where $\deg$ denotes total
degree. Write $m=\deg P$. If $\mathcal{E}^{\mathrm{alg}}_p(P)=0$, then
\[
  \mathcal{E}^{\mathrm{alg}}_p\bigl(H_m(P)\bigr)=0,
\]
where $H_m(P)$ is the $m$-th homogeneous component of $P$.
\end{lemma}
\begin{proof}
Put $m=\deg P\ge2$, $P_m=H_m(P)$ and $r=P-P_m$.

\emph{Step 1: $P_m$ is homogeneous of degree $m$ and $\deg r\le m-1$.} By the
background proposition on homogeneous components, $P_m=H_m(P)$ is homogeneous of
degree $m$ and $P=\sum_{d=0}^{m}H_d(P)$. Subtracting the term $d=m$,
\[
  r=P-P_m=\sum_{d=0}^{m-1}H_d(P).
\]
For each $0\le d\le m-1$ the polynomial $H_d(P)$ is homogeneous of degree $d$,
so every exponent tuple in its support has entry sum $d$ and hence
$\deg H_d(P)\le d\le m-1$. By the background proposition on the total degree of
a finite sum, $\deg r\le m-1$.

\emph{Step 2: a bookkeeping lemma.} We claim: if $\varphi\in\mathbb{R}[X_1,\dots,X_n]$
and $D\in\mathbb{N}$ satisfy $\deg\varphi\le D$, then
\begin{enumerate}
\item[(i)] for every index $a$, either $\partial_a\varphi=0$ or
  $\deg(\partial_a\varphi)+1\le D$;
\item[(ii)] for all indices $a,b$, either $\partial_a\partial_b\varphi=0$ or
  $\deg(\partial_a\partial_b\varphi)+2\le D$.
\end{enumerate}

For (i), suppose $\partial_a\varphi\ne0$. If $D=0$ then $\deg\varphi=0$, so
$\varphi$ is constant by the background proposition on total degrees, and then
$\partial_a\varphi=0$ by the background proposition on partial derivatives — a
contradiction. Hence $D\ge1$. By \noderef{TotalDegreePDeriv},
$\deg(\partial_a\varphi)\le\deg\varphi-1\le D-1$, and since $D\ge1$ this is the
same as $\deg(\partial_a\varphi)+1\le D$.

For (ii), suppose $\partial_a\partial_b\varphi\ne0$. Then $\partial_b\varphi\ne0$,
since $\partial_a0=0$. Put $D'=\deg(\partial_b\varphi)$. By (i) applied to
$\varphi$ and $D$ we get $D'+1\le D$. By (i) applied to the polynomial
$\partial_b\varphi$ with the bound $D'$ and the index $a$ (legitimate, since
$\partial_a(\partial_b\varphi)\ne0$) we get
$\deg(\partial_a\partial_b\varphi)+1\le D'$. Adding $1$ to both sides and
combining with $D'+1\le D$,
\[
  \deg(\partial_a\partial_b\varphi)+2\le D'+1\le D .
\]

\emph{Step 3: degree budgets for the factors.} Applying Step 2 to
$\varphi=P_m$ with $D=m$ (legitimate since $\deg P_m\le m$), and to
$\varphi=r$ with $D=m-1$ (legitimate by Step 1), we obtain, for all indices
$a,b$: every nonzero factor among
\[
  \partial_aP_m,\quad \partial_a\partial_bP_m,\quad \partial_ar,\quad
  \partial_a\partial_br
\]
satisfies
\begin{equation}\label{eq:tc-budget}
  \deg F+k\le m\ \ \text{if }F\text{ is a $k$-th derivative of }P_m,
  \qquad
  \deg F+k+1\le m\ \ \text{if }F\text{ is a $k$-th derivative of }r,
\end{equation}
where $k\in\{1,2\}$ is the order of the derivative. Indeed, for $P_m$ the two
cases are exactly (i) and (ii) with $D=m$; and for $r$ they are (i) and (ii)
with $D=m-1$, after adding $1$ to both sides and using $m-1+1=m$ (valid since
$m\ge2\ge1$).

\emph{Step 4: the expansion.} By the background proposition on additivity of
$\partial_i$, and since $P=P_m+r$,
\[
  \partial_aP=\partial_aP_m+\partial_ar,
  \qquad
  \partial_a\partial_bP=\partial_a\partial_bP_m+\partial_a\partial_br .
\]
Abbreviate $A_a=\partial_aP_m$, $\alpha_a=\partial_ar$,
$B_{ab}=\partial_a\partial_bP_m$, $\beta_{ab}=\partial_a\partial_br$. Expanding
the product of the two sums in the definition in \noderef{EpAlg} applied to $P$,
\[
  \mathcal{E}^{\mathrm{alg}}_p(P)
  =\sum_{i}\sum_{k}(A_i+\alpha_i)(A_i+\alpha_i)(B_{kk}+\beta_{kk})
   +(p-2)\sum_{i}\sum_{j}(B_{ij}+\beta_{ij})(A_i+\alpha_i)(A_j+\alpha_j).
\]
Expanding each triple product by distributivity gives $2^3=8$ summands. Exactly
one summand of each triple product has all three factors taken from $P_m$,
namely $A_iA_iB_{kk}$ in the first block and $B_{ij}A_iA_j$ in the second; and
\[
  \sum_i\sum_k A_iA_iB_{kk}=\Bigl(\sum_iA_i^2\Bigr)\Bigl(\sum_kB_{kk}\Bigr),
  \qquad
  \sum_i\sum_j B_{ij}A_iA_j,
\]
so these pure summands reassemble, again by the definition in \noderef{EpAlg},
into $\mathcal{E}^{\mathrm{alg}}_p(P_m)$. Collecting all the remaining summands
(seven from each triple product, over all index pairs, those from the second
block carrying the scalar $p-2$) into a single polynomial $M$, we have
\begin{equation}\label{eq:tc-split}
  \mathcal{E}^{\mathrm{alg}}_p(P)=\mathcal{E}^{\mathrm{alg}}_p(P_m)+M .
\end{equation}
By construction every summand of $M$ is a product of three factors, possibly
times the scalar $p-2$, in which exactly one factor is a second derivative
(order $2$) and two factors are first derivatives (order $1$), each factor is a
derivative of $P_m$ or of $r$, and \emph{at least one factor is a derivative of
$r$}.

\emph{Step 5: $\deg M\le 3m-5$.} Let $N$ be one of those summands, say
$N=c\,F_1F_2F_3$ with $c\in\{1,p-2\}$ and with $F_t$ a derivative of order
$k_t$, where $(k_1,k_2,k_3)$ is a permutation of $(1,1,2)$, so that
$k_1+k_2+k_3=4$.

If some $F_t=0$ then $N=0$ and $\deg N=0\le3m-5$, because $m\ge2$ gives
$3m-5\ge1$.

Otherwise all three factors are nonzero, so \eqref{eq:tc-budget} applies to each
of them: $\deg F_t+k_t\le m$ for every $t$, and moreover
$\deg F_t+k_t+1\le m$ for at least one $t$, namely any $t$ for which $F_t$ is a
derivative of $r$ (such a $t$ exists by Step 4). Adding the three inequalities,
taking the improved one for that distinguished $t$,
\[
  \bigl(\deg F_1+\deg F_2+\deg F_3\bigr)+(k_1+k_2+k_3)+1\le 3m,
\]
that is $\deg F_1+\deg F_2+\deg F_3+5\le3m$, i.e.
$\deg F_1+\deg F_2+\deg F_3\le3m-5$. By the background proposition on the total
degree of a product and of a scalar multiple,
\[
  \deg N\le\deg F_1+\deg F_2+\deg F_3\le3m-5 .
\]
Since $M$ is a finite sum of such summands, the background proposition on the
total degree of a finite sum gives $\deg M\le3m-5$.

\emph{Step 6: conclusion.} The operator $H_{3m-4}$ is $\mathbb{R}$-linear
(background proposition on homogeneous components). Applying it to
\eqref{eq:tc-split},
\[
  H_{3m-4}\bigl(\mathcal{E}^{\mathrm{alg}}_p(P)\bigr)
  =H_{3m-4}\bigl(\mathcal{E}^{\mathrm{alg}}_p(P_m)\bigr)+H_{3m-4}(M).
\]
We evaluate the three terms.
\begin{itemize}
\item The left-hand side is $H_{3m-4}(0)=0$, by the hypothesis
  $\mathcal{E}^{\mathrm{alg}}_p(P)=0$ and linearity.
\item By \noderef{EpAlgHomogeneous}, whose hypotheses $m\ge2$ and ``$P_m$ is
  homogeneous of degree $m$'' hold by assumption and by Step 1, the polynomial
  $\mathcal{E}^{\mathrm{alg}}_p(P_m)$ is homogeneous of degree $3m-4$. Hence
  $H_{3m-4}\bigl(\mathcal{E}^{\mathrm{alg}}_p(P_m)\bigr)=\mathcal{E}^{\mathrm{alg}}_p(P_m)$,
  since $H_d$ is the identity on polynomials homogeneous of degree $d$.
\item Since $m\ge2$ we have $3m-4\ge2$, so $3m-5=(3m-4)-1<3m-4$; combined with
  Step 5 this gives $\deg M<3m-4$, whence $H_{3m-4}(M)=0$, since $H_d$ kills
  polynomials of total degree $<d$.
\end{itemize}
Substituting these three evaluations,
$0=\mathcal{E}^{\mathrm{alg}}_p(P_m)+0$, i.e.
$\mathcal{E}^{\mathrm{alg}}_p\bigl(H_m(P)\bigr)=0$, which is the assertion.
\end{proof}
